\chapter{Описание пространственного движения}

Винтовое исчисление --- это раздел механики, изучающий винты. Винт в свою очередь это объект, хранящий в себе информацию о векторной и моментной частях движения в трехмерном пространстве. Основы теории были заложены в конце XIX века, наиболее подробно теори винтов изложена в трудах Ф.М. Диментберга ~\cite{Dimentberg:1965}, Ю.Н. Челнокова ~\cite{Chelnokov:2006} и В.Н. Гордеева ~\cite{Gordeev:2016}.
\section{Дуальные числа}
Дуальные числа --- это числа вида:
\begin{equation*}
    U = u + \varepsilon u^o, \varepsilon^2 = 0,
\end{equation*}
где $u, u^\circ \in \mathbb{R}$, а $\varepsilon$ является символом Клиффорда для которого выполняется определяющее свойство $\varepsilon^2 = 0$ при $\varepsilon \neq 0$. В этой записи $u$ и $u^\circ$ называются соответственно $\emph{главной}$ и $\emph{моментной}$ частями. Благодаря свойствам символа $\varepsilon$, дуальные числа также называют $\emph{параболическими числами Клиффорда}$.
Важно заметить, что при $u \neq 0$ число представляется через параметр $p = \frac{u^o}{u}$:
\begin{equation*}
    U = u(1 + \varepsilon u) = u e^{\varepsilon p}.
\end{equation*}

Как и с любыми другими числами, над дуальными числами можно проводить операции, такие как:
\begin{enumerate}
    \item $U \pm K = (u \pm k) + \varepsilon (u^o \pm k^o)$ --- сложение и вычитание,
    \item $UK = uk + \varepsilon (u k^o + u^o k)$ --- умножение,
    \item $\frac{U}{K} = \frac{u}{k} + \varepsilon \frac{u^o k - u k^o}{k^2}$ --- деление с учетом $k \neq 0$,
    \item $U^n = u^n + \varepsilon n u^o u^{n-1}, \quad U^{1/n} = u^{1/n} + \varepsilon \frac{u^o}{n} u^{1/n - 1}$ --- возведение в степень и извлечение корня.
\end{enumerate}

Также нельзя делить число с ненулевой главной частью на число с нулевой главной частью, потому что это приведет к обнулению выражения, аналогично и для произведения.

В дальнейшем мы будем рассматривать винтовые преобразования, где дуальные числа понадобятся для описания движения твердого тела.

Если $U = u + \varepsilon u^o, \varepsilon^2 = 0$, то разложение любой аналитической функции $F(U)$ имеет следующий вид:
\begin{equation*}
    F(U) = F(u + \varepsilon u^o) = f(u) + \varepsilon u^o \frac{df(u)}{du},
\end{equation*}
где где $f(u) = F(u)$ --- главная часть функции.

Также, если функция зависима от дуальных параметров $K_1 = k_1 + \varepsilon k_1^o, \quad \dots, \quad K_n = k_n + \varepsilon k_n^o$, то;
\begin{equation*}
    F(U, K_1, \dots, K_n) = f(u) + \varepsilon \big[ u^o f'(u) + f^o(u) \big],
\end{equation*}
где 
\begin{equation*}
    f(u) = F(u, k_1, \dots, k_n), \quad f'(u) = \frac{\partial f}{\partial u},
\end{equation*}
\begin{equation*}
    f^o(u) = k_1^o \frac{\partial f}{\partial k_1} + \dots + k_n^o \frac{\partial f}{\partial k_n}.
\end{equation*}

Рассмотрим примеры элементарных функций:
\begin{equation*}
    \sin U = \sin u + \varepsilon u^o \cos u,
\end{equation*}
\begin{equation*}
    \cos U = \cos u - \varepsilon u^o \sin u,
\end{equation*}
\begin{equation*}
    \ln U = \ln u + \varepsilon \frac{u^o}{u}, \quad u > 0,
\end{equation*}
\begin{equation*}
    e^{U} = e^{u}(1 + \varepsilon u^o).
\end{equation*}

Для функции дуальной переменной производная будет иметь следующий вид:
\begin{equation*}
    \frac{dF(U)}{dU} = f'(u) + \varepsilon \big[ u^o f''(u) + f^{o'}(u) \big].
\end{equation*}
Для интеграла будет верно:
\begin{equation*}
    \int F(U) \, dU = \int f(u) \, du + \varepsilon \big[ u^o f(u) + \int f^o(u) \, du \big] + C,
\end{equation*}
где $C = c + \varepsilon c^o$ --- дуальная постоянная.

В случае и с интегрированием, и с дифференцированием с дуальной переменной можно свести к обычному:
\begin{equation*}
    \frac{dF(U)}{dU} = \frac{\partial F(U)}{\partial u}, \quad \int F(U) \, dU \equiv \int F(U) \, du.
\end{equation*}

Смысл в том, что все привычные прежде формулы остаются верными и для дуальных переменных. Их функции определяются главной частью $f(u)$, а дуальная часть получается в результате производной от главной части по главному аргументу и если есть, то через параметры.

\section{Скользящие векторы}

В евклидовом пространстве выделяют три вида векторов:
\begin{enumerate}
    \item Свободный вектор --- характеризуется только направлением и длиной.
    \item Скользящий вектор --- привязан к прямой и может совершать перемешение только вдоль этой прямой.
    \item Неподвижный вектор --- привязан к конкретной точке приложения.
\end{enumerate}

В контексте скользящего вектора, также упоминаются координаты Плюккера. Скользящий вектор может быть задан шестеркой чисел в трехмерном пространстве:
\begin{equation*}
    \vb{R} = \vb{r} + \varepsilon \vb{r}^o = \begin{bmatrix} r_1 + \varepsilon r_1^o \\ r_2 + \varepsilon r_2^o \\ r_3 + \varepsilon r_3^o \end{bmatrix},
\end{equation*}
где $\vb{r} = (r_1, r_2, r_3)^T$ является главной частью, отвечающей за направление и длину, а $\vb{r}^o = (r_1^o, r_2^o, r_3^o)^T$ моментной частью(момент относительно начала координат).

Главная и моментная части связаны, поэтому моментную часть можно выразить через радиус-вектор $\vb{\rho}$ любой точки на линии вектора:
\begin{equation*}
    r_1^o = \rho_2 r_3 - r_2 \rho_3, \quad r_2^o = \rho_3 r_1 - r_3 \rho_1, \quad r_3^o = \rho_1 r_2 - r_1 \rho_2,
\end{equation*}
а в векторной форме:
\begin{equation*}
    \vb{r}^o = \vb{\rho} \times \vb{r}.
\end{equation*}

Существует условие Плюккера, которое заключается в том, что для одного скользящего вектора будет выполнять условие ортогональности:
\begin{equation*}
    r_1 r_1^o + r_2 r_2^o + r_3 r_3^o = 0 \quad \Leftrightarrow \quad \vb{r} \cdot \vb{r}^o = 0.
\end{equation*}

По координатам Плюккера также можно найти точку на линии вектора, которая будет являться ближайшей к началу координат:
\begin{equation*}
    d_1 = \frac{r_2 r_3^o - r_2^o r_3}{r_1^2 + r_2^2 + r_3^2}, \quad d_2 = \frac{r_3 r_1^o - r_3^o r_1}{r_1^2 + r_2^2 + r_3^2}, \quad d_3 = \frac{r_1 r_2^o - r_1^o r_2}{r_1^2 + r_2^2 + r_3^2}.
\end{equation*}
В векторной форме:
\begin{equation*}
    \vb{d} = \frac{\vb{r} \times \vb{r}^o}{\|\vb{r}\|^2}.
\end{equation*}
У найденной точки будет ряд свойств:
\begin{itemize}
    \item Радиус вектор $\vb{d}$ перпендикулярен главному вектору $\vb{r}$.
    \item Векторы $\vb{r}$, $\vb{r}^o$ и $\vb{d}$ попарно ортогональны.
\end{itemize}

Сложение плюккеровых координат является их преимуществом так, как выполняется покомпонентно:
\begin{equation*}
    \vb{R}_\Sigma = \vb{r}_\Sigma + \varepsilon \vb{r}_\Sigma^o,
\end{equation*}
где при сложении главный вектор $\vb{r}_\Sigma = \sum \vb{r}_k$, а главный момент $\vb{r}_\Sigma^o = \sum \vb{r}_k^o$.

Отметим, что после сложение условие Плюккера может нарушиться:
\begin{equation*}
    \vb{r}_\Sigma \cdot \vb{r}_\Sigma^o \neq 0.
\end{equation*}

Введем также такие понятия, как инвариант дуального вектора:
\begin{equation*}
    I = \vb{r}_\Sigma \cdot \vb{r}_\Sigma^o = r_1 r_1^o + r_2 r_2^o + r_3 r_3^o.
\end{equation*}
и параметр дуального вектора:
\begin{equation*}
    [\vb{R}] = \frac{\vb{r}_\Sigma \cdot \vb{r}_\Sigma^o}{\|\vb{r}_\Sigma\|^2}.
\end{equation*}

\section{Винты}

Винт, а в терминологии Котельникова упоминается как мотор, явялется математическим объектом, который может иметь представление в виде дуального вектора:
\begin{equation*}
    \vb{R} = \vb{r} + \varepsilon \vb{r}^o = \begin{bmatrix} r_1 + \varepsilon r_1^o \\ r_2 + \varepsilon r_2^o \\ r_3 + \varepsilon r_3^o \end{bmatrix},
\end{equation*}
где $\varepsilon^2 = 0$.

К частным случаям относят скользящий вектор, винт и пару сил. Отметим, что в случае с скользящим вектором момент ортогонален вектору, у винта момент будет коллинеарен вектору, а у пары сил есть чистый момент, который не привязан к вектору силы .

В теории винтов вводятся такие понятия, как норма винта или как ее еще называют дуальная величина:
\begin{equation*}
    \|\vb{R}\| = (\vb{r} + \varepsilon \vb{r}^o) \cdot (\vb{r} + \varepsilon \vb{r}^o) = \|\vb{r}\|^2 + 2\varepsilon (\vb{r} \cdot \vb{r}^o).
\end{equation*}
Модуль винта(дуальное число):
\begin{equation*}
    |\vb{R}| = \sqrt{\|\vb{R}\|} = \|\vb{r}\| \left(1 + \varepsilon \frac{\vb{r} \cdot \vb{r}^o}{\|\vb{r}\|^2}\right) = \|\vb{r}\| (1 + \varepsilon [\vb{R}]),
\end{equation*}
где $[\vb{R}] = \frac{\vb{r} \cdot \vb{r}^o}{\|\vb{r}\|^2}$ является параметром винта.Заметим, что понятие модуля определяется только при $\vb{r} \neq \vb{0}$.

Единичный винт представляет собой ось с направлением в пространстве и обозначается, как $\vb{E}$. Для него будет верна следующая запись:
\begin{equation*}
    \vb{E} = \vb{e} + \varepsilon \vb{e}^o, 
\end{equation*}
где $e_1^2 + e_2^2 + e_3^2 = 1$, $e_1 e_1^o + e_2 e_2^o + e_3 e_3^o = 0$.Такой винт обладает единичной векторной частью $\|\vb{e}\| = 1$ и параметром равным нулю $\vb{e} \cdot \vb{e}^o = 0$.

Винт имеет два основных представления:
\begin{enumerate}
    \item Через плюккеровы координаты или бивекторное:
        \begin{equation*}
            \vb{R} = \vb{r} + \varepsilon \vb{r}^o.
        \end{equation*}
    \item Через ось и дуальный модуль(винтовое представление):
        \begin{equation*}
            \vb{R} = R \vb{E},
        \end{equation*}
        где $R = |\vb{R}| = \|\vb{r}\| (1 + \varepsilon [\vb{R}])$.
\end{enumerate}

\subsection{Операции над винтами}

Сложение винтов выполняется покомпонентно:
\begin{equation*}
    \vb{R} = \vb{R}_1 + \vb{R}_2 + \dots + \vb{R}_n,
\end{equation*}
при условиях $\vb{r} = \vb{r}_1 + \vb{r}_2 + \dots + \vb{r}_n$ и $\vb{r}^o = \vb{r}_1^o + \vb{r}_2^o + \dots + \vb{r}_n^o$.

Сложение винтов является коммутативным и ассоциативным. Скалярное произведение дистрибутивно:
\begin{equation*}
    \vb{R} \cdot \vb{S} = (\vb{R}_1 + \vb{R}_2 + \dots) \cdot \vb{S} = \vb{R}_1 \cdot \vb{S} + \vb{R}_2 \cdot \vb{S} + \dots.
\end{equation*}

Умножение винта может быть произведено на разные виды чисел. Например, для умножения на действительное число будет верно:
\begin{equation*}
    a \vb{R} = a \vb{r} + \varepsilon a \vb{r}^o,
\end{equation*}
здесь сохраняется ось винта и при $a > 0$ направление сохраняется, при $a < 0$ меняется на противоположное.

Для умножения на дуальное число вида $A = a + \varepsilon a^o$:
\begin{equation*}
    A \vb{R} = (a + \varepsilon a^o)(\vb{r} + \varepsilon \vb{r}^o) = a \vb{r} + \varepsilon (a \vb{r}^o + a^o \vb{r}),
\end{equation*}
а в винтовом представлении примет следующий вид:
\begin{equation*}
    A \vb{R} = (a + \varepsilon a^o) \|\vb{r}\| e^{\varepsilon [\vb{R}]} \vb{E} = a \|\vb{r}\| e^{\varepsilon ([\vb{R}] + a^o / a)} \vb{E}.
\end{equation*}

Перейдем к скалярному произведению винтов. Для двух винтов вида $\vb{R}_1 = \vb{r}_1 + \varepsilon \vb{r}_1^o$ и $\vb{R}_2 = \vb{r}_2 + \varepsilon \vb{r}_2^o$ будет верна следующая запись:
\begin{equation*}
    \vb{R}_1 \cdot \vb{R}_2 = \vb{r}_1 \cdot \vb{r}_2 + \varepsilon (\vb{r}_1 \cdot \vb{r}_2^o + \vb{r}_1^o \cdot \vb{r}_2).
\end{equation*}
В винтовом представлении:
\begin{equation*}
    \vb{R}_1 \cdot \vb{R}_2 = R_1 R_2 \cos \Gamma,
\end{equation*}
для $\vb{R}_1 = R_1 \vb{E}_1$, $\vb{R}_2 = R_2 \vb{E}_2$ и дуальный угол между осями $\Gamma = \gamma + \varepsilon \gamma^o$.Винтовое произведение также является коммутативным $\vb{R}_1 \cdot \vb{R}_2 = \vb{R}_2 \cdot \vb{R}_1$.

Для винтового же произведения будет верно:
\begin{equation*}
    \vb{R}_1 \times \vb{R}_2 = \vb{r}_1 \times \vb{r}_2 + \varepsilon (\vb{r}_1 \times \vb{r}_2^o + \vb{r}_1^o \times \vb{r}_2),
\end{equation*}
а винтовое представление примет следующий вид:
\begin{equation*}
    \vb{R}_1 \times \vb{R}_2 = R_1 R_2 \sin \Gamma \, \vb{E}_{12},
\end{equation*}
где $\vb{E}_{12}$ является единичным винтом, направленным по общему перпендикуляру.Здесь нет коммутативности, как в случае со скалярным произведением, винтовое произведение обладает свойством антикоммутативности $\vb{R}_1 \times \vb{R}_2 = -\vb{R}_2 \times \vb{R}_1$. Также ось результирующего винта (произведения) перпендикулярна осям исходных и модуль равен произведению дуальных модулей на синус дуального угла.

К операциям над винтами также относят проекцию винта на ось. Для проекции винта $\vb{R}$ на ось, заданную единичным винтом $\vb{E}$ будет выполняться:
\begin{equation*}
    \vb{R} \cdot \vb{E} = r \cos \gamma + \varepsilon r (p \cos \gamma - \gamma^o \sin \gamma)
\end{equation*}

\subsection{Дуальный угол между двумя осями}

Рассмотрим две нправленные оси $O_1$ и $O_2$, тогда дуальный угол между ними:
\begin{equation*}
    \Gamma = \gamma + \varepsilon \gamma^o,
\end{equation*}
где $\gamma$ является обычным углом между осями $(0 \leq \gamma \leq \pi)$, а $\gamma^o$ кратчайшее расстояние между осями(знак определяется по правилу правой тройки).

Для двух единичных винтов $\vb{O}_1 = \vb{o}_1 + \varepsilon \vb{o}_1^o$ и $\vb{O}_2 = \vb{o}_2 + \varepsilon \vb{o}_2^o$ будет верно:
\begin{equation*}
    \cos \Gamma = \vb{O}_1 \cdot \vb{O}_2 = \vb{o}_1 \cdot \vb{o}_2 + \varepsilon (\vb{o}_1 \cdot \vb{o}_2^o + \vb{o}_1^o \cdot \vb{o}_2),
\end{equation*}
а координатное выражение примет вид:
\begin{equation*}
    \cos \gamma = o_{11} o_{21} + o_{12} o_{22} + o_{13} o_{23},
\end{equation*}
\begin{equation*}
    \gamma^o \sin \gamma = o_{11} o_{21}^o + o_{12} o_{22}^o + o_{13} o_{23}^o + o_{11}^o o_{21} + o_{12}^o o_{22} + o_{13}^o o_{23}.
\end{equation*}

Дуальный угол также обладает рядом тригонометрических свойств:
\begin{itemize}
    \item $\sin \Gamma = \sin \gamma + \varepsilon \gamma^o \cos \gamma$,
    \item $\cos \Gamma = \cos \gamma - \varepsilon \gamma^o \sin \gamma$,
    \item $\tg \Gamma = \tg \gamma + \frac{\varepsilon \gamma^o}{\cos^2 \gamma}$.
\end{itemize}

Рассмотрев основные тригонометрические свойства можно перейти к координатному представлению винтов. Для описания такого представления используются дуальные координаты. Винт $\vb{R}$ можно представить через три дуальные координаты, а если точнее проекции на оси прямоугольной системы координат. Соответственно, для $X$;
\begin{equation*}
    X = x + \varepsilon x^o = R \cos \Gamma_X = r[\cos \gamma_X + \varepsilon (p \cos \gamma_X - \gamma_X^o \sin \gamma_X)].
\end{equation*}
Для $Y$:
\begin{equation*}
    Y = y + \varepsilon y^o = R \cos \Gamma_Y = r[\cos \gamma_Y + \varepsilon (p \cos \gamma_Y - \gamma_Y^o \sin \gamma_Y)].
\end{equation*}
Для $Z$:
\begin{equation*}
    Z = z + \varepsilon z^o = R \cos \Gamma_Z = r[\cos \gamma_Z + \varepsilon (p \cos \gamma_Z - \gamma_Z^o \sin \gamma_Z)].
\end{equation*}
Здесь $\Gamma_X, \Gamma_Y, \Gamma_Z$ являются дуальными углами между осью винта и координатными осями.

Винт также возможно представить в виде винта:
\begin{equation*}
    \vb{R} = X \vb{i} + Y \vb{j} + Z \vb{k} = (x + \varepsilon x^o)\vb{i} + (y + \varepsilon y^o)\vb{j} + (z + \varepsilon z^o)\vb{k}.
\end{equation*}

К основным свойствам подобного представления относят:
\begin{itemize}
    \item Квадрат модуля:
        \begin{equation*}
            R^2 = X^2 + Y^2 + Z^2 = (x^2 + y^2 + z^2) + 2\varepsilon (x x^o + y y^o + z z^o).
        \end{equation*}
    \item Основное тождество для дуальных направляющих косинусов:
        \begin{equation*}
            \cos^2 \Gamma_X + \cos^2 \Gamma_Y + \cos^2 \Gamma_Z = 1,
        \end{equation*}
        впоследствии это тождество распадается на два вещественных уравнения $\cos^2 \gamma_X + \cos^2 \gamma_Y + \cos^2 \gamma_Z = 1$ и $\gamma_X^o \cos \gamma_X \sin \gamma_X + \gamma_Y^o \cos \gamma_Y \sin \gamma_Y + \gamma_Z^o \cos \gamma_Z \sin \gamma_Z = 0$.
\end{itemize}
Для скалярного произведения в координатах будет верно:
\begin{equation*}
    \vb{R}_1 \cdot \vb{R}_2 = X_1 X_2 + Y_1 Y_2 + Z_1 Z_2 = (x_1 x_2 + y_1 y_2 + z_1 z_2) + \varepsilon (x_1 x_2^o + x_1^o x_2 + y_1 y_2^o + y_1^o y_2 + z_1 z_2^o + z_1^o z_2),
\end{equation*}
а для винтового:
\begin{equation*}
    \vb{R}_1 \times \vb{R}_2 = \begin{vmatrix} \vb{i} & \vb{j} & \vb{k} \\ X_1 & Y_1 & Z_1 \\ X_2 & Y_2 & Z_2 \end{vmatrix}.
\end{equation*}

\section{Пространственные перемещения твердого тела с помощью кинематических параметров}

Мы будем рассматривать движение свободного твердого тела относительно произвольной неподвижной системы координат $O X_1 X_2 X_3$. Основное отличие от тела с закрепеленной точкой в том, что здесь тело можемт как вращаться, так и перемещаться поступательно.

Согласно теореме Шаля ~\cite{Chelnokov:2006}:
\begin{quote}
    <<Всякое перемещение свободного твердого тела из одного положения в другое может быть получено посредством поступательного перемещения тела вместе с произвольно выбранным полюсом и поворота тела вокруг некоторой оси, проходящей через этот полюс.>>
\end{quote}
Следовательно поступательная часть зависит от выбора полюса, так как разные полюса дают разные векторы переноса, а вращательная часть в свою очередь не зависит от выбора.

Для бесконечно малых перемещений теорема Шаля подводит к понятию винтового движения. Винтовое движение --- это комбинаяция вращения вокруг некоторой оси с угловой скоростью $\omega$ и поступательного движения вдоль той же оси со скоростью $v$. Параметр винта(шаг) имеет следующий вид:
\begin{equation*}
    h = \frac{v}{\omega},
\end{equation*}
при $\omega \neq 0$. Если же $\omega = 0$, то движение будет чисто поступательным, если $v = 0$ то чисто вращательное.

При движении твердого тела мгновенная винтовая ось(ось, вокруг который происходит вращение и вдоль которой поступательное движение) меняет свое положение в пространстве.
Неподвижный аксоид --- это место мгновенных винтовых осей в неподвижной системе координат, а подвижный --- место мгновенных винтовых осей в системе координат, связанной с телом.

Оба аксоида представляют собой линейчатые поверхности, которые не развертываются, т.е. образованы движением прямой линии. В частных случаях аксоиды могут вырождаться в конические поверхности при движении с закрепленной точкой и цилиндрические поверхности при плоском движении. 

К традиционным конематическим параметрам, которые используют для описания движения твердого тела, относят:
\begin{itemize}
    \item Поступательные параметры --- декартовы координаты выбранного полюса, всего 3 параметра.
    \item Вращательные параметры --- углы, задающие ориентацию тела, также всего 3 параметра.
\end{itemize}

\subsection{Дуальные углы}

Для задания ориентации твердого тела используются углы Эйлера $\psi, \theta, \varphi$ и углы Крылова, их еще называют самолетными.

Для описания свободного твердого тела, которое может и вращаться, и поступательно перемещаться, логичнее всего использовать дуальные аналоги обычных углов --- дуальные углы Эйлера-Крылова.
Дуальный угол же это объект вида:
\begin{equation*}
    \Theta = \theta + \varepsilon \theta^o,
\end{equation*}
где $\theta$ вляется обычным углом поворота, $\theta^o$ линейным перемещением вдоль оси вращения, $\varepsilon$ дуальной единицей.

Если же рассматривать отдельно дуалные углы Эйлера, то здесь важно отметить, что происходит переход из начального положения в произвольное конечное с помощью последовательности трех винтовых перемещений, которые как раз и задаются этими углами.
\begin{itemize}
    \item Дуальный угол прецессии:
        \begin{equation*}
            \Psi = \psi + \varepsilon \psi^o,
        \end{equation*}
        где вращение происходит вокруг оси $X_3$ или $Y_3$.
    \item Дуальный угол нутации:
        \begin{equation*}
            \Theta = \theta + \varepsilon \theta^o,
        \end{equation*}
        где вращение происходит вокруг новой оси $Y_1'$.
    \item Дуальный угол собственного вращения:
        \begin{equation*}
            \Phi = \varphi + \varepsilon \varphi^o,
        \end{equation*}
        где вращение происходит вокруг конечной оси $Y_3$.
\end{itemize}

Важно не забывать о правилах знаков при рассмотрении подобных пространственных перемещений. Углы $\psi, \theta, \varphi$ положительны при вращении против часовой стрелки, если смотреть относительно конца соответствующей оси. Линейные перемещения  $\psi^o, \theta^o, \varphi^o$ положительны в направлении соответствующей оси.

Говоря о недостатках таких углов можно заметить аналогию с обычными углами Эйлера, здесь  присуствует отсутствие симметрии. При малых отклонениях от начального положения малыми будут только $\theta$ и сумма $\psi + \varphi$, также отдельные углы $\psi$ и $\varphi$ могут принимать произвольные значения. Это все приводит к сингулярностям, например, при $\theta=0$, когда оси $X_3$ и $Y_3$ совпадают, собственное вращение становится неразличимым.

Как и ранее упоминалось, существуют самолетные углы --- дуальные угла Крылова. Они являются альтернативой, учитывая все недостатки углов Эйлера. Они называются самолетными так как получили применение в авиации, поэтому будем использоваться соответствующие наименования:
\begin{itemize}
    \item Дуальный угол рыскания:
        \begin{equation*}
            \Psi = \psi + \varepsilon \psi^o,
        \end{equation*}
        где вращение происходит вокруг нормальной оси $X_2$.
    \item Дуальный угол тангажа:
        \begin{equation*}
            \Theta = \theta + \varepsilon \theta^o,
        \end{equation*}
        где вращение происходит вокруг поперечной оси.
    \item Дуальный угол крена:
        \begin{equation*}
            \Phi = \varphi + \varepsilon \varphi^o,
        \end{equation*}
        где вращение происходит вокруг продольной оси.
\end{itemize}

В случае с углами Крылова отсутсвует проблема сингулярностей и при малых отклонения от начального положения малыми будут все три дуальных угла.

\subsection{Дуальные направляющие косинусы}

Для описания ориентации твердого тела вместо углов можно использовать направляющие косинусы. В дуальном варианте это позволяет одновременно описывать и ориентацию, и положение. Рассмотрим две системы координат, $X$ --- неподвиженая система $O X_1 X_2 X_3$ и $Y$ --- система, связанная с телом $A Y_1 Y_2 Y_3$.
Дуальный направляющий косинус --- это скалярное произведение единичных векторов осей:
\begin{equation*}
    C_{ik} = \vb{Y}_i \cdot \vb{X}_k = \cos \Gamma_{ik},
\end{equation*}
где дуальныи углом между осями $Y_i$ и $X_k$ является $\Gamma_{ik} = \gamma_{ik} + \varepsilon \gamma_{ik}^o$.

Если разложить дуальный косинус:
\begin{equation*}
    C_{ik} = \underbrace{\cos \gamma_{ik}}_{c_{ik}} + \varepsilon \underbrace{(-\gamma_{ik}^o \sin \gamma_{ik})}_{c_{ik}^o},
\end{equation*}
то можно увидеть, что $\gamma_{ik}$ обычный угол между осями и $\gamma_{ik}^o$ кратчайшее расстояние между осями.

Матрица для дуальный направляющих косинусов имеет следующий вид:
\begin{equation*}
    \vb{C} = \|C_{ik}\| = \|c_{ik} + \varepsilon c_{ik}^o\| = \begin{pmatrix}
        C_{11} & C_{12} & C_{13} \\
        C_{21} & C_{22} & C_{23} \\
        C_{31} & C_{32} & C_{33}
        \end{pmatrix}
\end{equation*}

К тому же, у Лурье ~\cite{Lurie:1961} уже были выведены формулы для элементов матрицы поворота через параметры Родрига-Гамильтона:
\begin{equation*}
    a_{sk} = \delta_{sk}\left(1 - 2\sum_{q=1}^3 \lambda_q^2\right) + 2\lambda_0 \sum_{t=1}^3 \varepsilon_{tsk}\lambda_t + 2\lambda_s\lambda_k.
\end{equation*}

\subsection{Винт конечного перемещения твердого тела}

Согласно теореме Шаля, любое перемещение сводобного твердого тела можно представить как винтовое перемещение, т.е. вращение вокруг некоторой оси и поступательное движение вдоль той же оси. Конечное винтовое перемещение описывается дуальным вектором:
\begin{equation*}
    \vb{\Theta} = 2\vb{E} \tg \frac{\Phi}{2},
\end{equation*}
где $\vb{E} = \vb{e} + \varepsilon \vb{e}^o$ является единичным винтом оси винтового перемещения, $\Phi = \varphi + \varepsilon \varphi^o$ дуальным углом поворота, $\varphi$ углом поворота вокруг оси, а $\varphi^o$ поступательным перемещением вдоль оси.

Проекция же винта $\vb{\Theta}$ на оси неподвижной системы $X$ имеет следующий вид:
\begin{equation*}
    \Theta_i = 2 \tg \frac{\Phi}{2} \cos \Gamma_i, \quad i = 1, 2, 3,
\end{equation*}
где $\Gamma_i = \gamma_i + \varepsilon \gamma_i^o$ --- дуальный угол между осью винта и осью $X_i$.

При рассмотрении конечных перемещений в пространстве также используются дуальные параметры Эйлера-Родрига-Гамильтона и Кейли-Клейна. У Лурье ~\cite{Lurie:1961}, например, параметры вводятся через вектор конечного поворота $\vb{\theta}$, что демонстрирует связь между параметрами Родрига-Гамильтона и компонентами вектора конечного поворота:
\begin{equation*}
    \frac{1}{2}\theta_s = \frac{\lambda_s}{\lambda_0}, \quad s = 1,2,3.
\end{equation*}

Дуальные параметры Эйлера, которые являются обобщением классических параметров Родрига-Гамильтона, определяются через проекции винта перемещения:
\begin{equation*}
    \Lambda_0 = \cos \frac{\Phi}{2}, \quad \Lambda_i = \frac{1}{2} \Lambda_0 \Theta_i = \sin \frac{\Phi}{2} \cos \Gamma_i, \quad i = 1, 2, 3
\end{equation*}



У дуальных параметров Эйлера–Родрига–Гамильтона каждый параметр является дуальным числом:
\begin{equation*}
    \Lambda_j = \lambda_j + \varepsilon \lambda_j^o, \quad j = 0, 1, 2, 3,
\end{equation*}
где $\lambda_j$ являются вещественными параметрами Эйлера, а $\lambda_j^o$ их дуалными компонентами.

Для дуальных параметров выполняется одно основное тождество:
\begin{equation*}
    \Lambda_0^2 + \Lambda_1^2 + \Lambda_2^2 + \Lambda_3^2 = 1.
\end{equation*}
Оно же распадается на два вещественных уравнения:
\begin{equation*}
    \lambda_0^2 + \lambda_1^2 + \lambda_2^2 + \lambda_3^2 = 1,
\end{equation*}
которое является условием нормировки главных частей и условие ортогональности:
\begin{equation*}
    \lambda_0 \lambda_0^o + \lambda_1 \lambda_1^o + \lambda_2 \lambda_2^o + \lambda_3 \lambda_3^o = 0.
\end{equation*}

Сами параметры Кейли-Клейна вводятся через параметры Эйлера:
\begin{equation*}
    \Alpha = \Lambda_0 + i\Lambda_3 = \alpha + \varepsilon \alpha^o,
\end{equation*}
\begin{equation*}
    \Beta = -\Lambda_2 + i\Lambda_1 = \beta + \varepsilon \beta^o,
\end{equation*}
\begin{equation*}
    \Gamma = \Lambda_2 + i\Lambda_1 = \gamma + \varepsilon \gamma^o,
\end{equation*}
\begin{equation*}
    \Delta = \Lambda_0 - i\Lambda_3 = \delta + \varepsilon \delta^o,
\end{equation*}
где $\alpha, \beta, \gamma, \delta$ вещественные параметры Кейли-Клейна, $\alpha^o, \beta^o, \gamma^o, \delta^o$ их дуальные компоненты, а $i$ мнимая единица.Основным их тождеством является $A\Delta - B\Gamma = 1$

\section{Формула Родрига и ее дуальный аналог}
\label{sec:Rodrig}

Классическая формула Родрига используетс для поворота точки вокруг оси, но также может быть обобщена на случай винтового движения. Сама формула Родрига в компактной векторной форме принимает следующий вид:
\begin{equation*}
    \vb{r}' = \vb{r} + \left(1 + \frac{\theta^2}{4}\right)^{-1} \left[ \vb{\theta} \times \left( \vb{r} + \frac{1}{2} \vb{\theta} \times \vb{r} \right) \right].
\end{equation*}
Но существуют и альтернативные формы записи, например, через тригонометрические функции:
\begin{equation*}
    \vb{r}' = \vb{r} + \sin \varphi (\vb{e} \times \vb{r}) + (1 - \cos \varphi) [\vb{e} \times (\vb{e} \times \vb{r})].
\end{equation*}
Или более развернутая:
\begin{equation*}
    \vb{r}' = \cos \varphi \, \vb{r} + \sin \varphi (\vb{e} \times \vb{r}) + (1 - \cos \varphi)(\vb{e} \cdot \vb{r})\vb{e}.
\end{equation*}

Существует теорема о двух полуоборотах согласно которой:
\begin{quote}
    <<Поворот твердого тела на угол \( \varphi \) вокруг некоторой оси эквивалентен двум последовательным полуоборотам тела вокруг осей, пересекающих под прямым углом в одной точке данную ось и образующих между собой угол $\varphi/2$.
\end{quote}

Для свободного твердого тела(нет закрепеленных точек) обобщением формулы Родрига является дуальный аналог, описывающий винтовое движение, как комбинацию вращательного и поступательного действий.
Введем дуальный аналоги величин, использованных в классической формуле:
\begin{itemize}
    \item Угол поворота $\varphi$ заменяется на дуальный угол $\Phi = \varphi + \varepsilon \varphi^\circ$.
    \item Единичный вектор оси $\vb{e}$ на единичный винт оси $\vb{E} = \vb{e} + \varepsilon \vb{e}^\circ$.
    \item Вектор конечного поворота $\vb{\theta} = 2\tg(\varphi/2)\vb{e}$ на дуальный вектор конечного перемещения $\vb{\Theta} = 2\vb{E} \tg(\Phi/2)$.
    \item Радиус-вектор точки $\vb{r}$ на винт точки $\vb{R} = \vb{r} + \varepsilon \vb{r}^\circ$.
\end{itemize}
В таком случае дуальная формула Родрига будет иметь следующий вид:
\begin{equation*}
    \vb{R}' = \vb{R} + \left(1 + \frac{\vb{\Theta}^2}{4}\right)^{-1} \left[ \vb{\Theta} \times \left( \vb{R} + \frac{1}{2} \vb{\Theta} \times \vb{R} \right) \right],
\end{equation*}
или в более развернутой тригонометрической форме:
\begin{equation*}
    \vb{R}' = \cos \Phi \, \vb{R} + \sin \Phi (\vb{E} \times \vb{R}) + (1 - \cos \Phi)(\vb{E} \cdot \vb{R})\vb{E}.
\end{equation*}

Существуют частные случаи перемещений в пространстве с использованием этой формулы.
\begin{itemize}
    \item Чистое вращение($\varphi^\circ = 0, \vb{e}^\circ = 0$):
        \begin{equation*}
            \vb{R}' = \vb{r}' + \varepsilon \vb{r}'^\circ.
        \end{equation*}
    \item Чистая трансляция($\varphi = 0, \vb{e}^\circ = 0$) и при $\varphi = 0$ , $\cos \Phi = 1$, $\sin \Phi = \varepsilon \varphi^\circ$:
        \begin{equation*}
            \vb{R}' = \vb{R} + \varepsilon \varphi^\circ (\vb{e} \times \vb{R}),
        \end{equation*}
        если с учетом главной и моментной частей:
        \begin{equation*}
            \vb{r}' = \vb{r}, \quad \vb{r}'^\circ = \vb{r}^\circ + \varphi^\circ (\vb{e} \times \vb{r}),
        \end{equation*}
        что полностью соответствует переносу точки на вектор $\varphi^\circ \vb{e}$, т.е. моментная часть изменяется.
    \item Полуоборот($\varphi = \pi, \varphi^\circ = 0$):
        \begin{equation*}
            \vb{r}' = 2(\vb{e} \cdot \vb{r})\vb{e} - \vb{r},
        \end{equation*}
        \begin{equation*}
            \vb{r}'^\circ = 2(\vb{e} \cdot \vb{r}^\circ + \vb{e}^\circ \cdot \vb{r})\vb{e} + 2(\vb{e} \cdot \vb{r})\vb{e}^\circ - \vb{r}^\circ.
        \end{equation*}
\end{itemize}

Ранее упоминалась теорема о двух полуоборотах, соответственно, у нее тоже есть дуальный аналог:
\begin{quote}
    <<Винтовое перемещение тела на дуальный угол $\Phi = \varphi + \varepsilon \varphi^\circ$ относительно оси, единичный винт которой $\vb{E}$, эквивалентно двум последовательным полуоборотам, совершаемым относительно прямых с единичными винтами $\vb{E}_1$ и $\vb{E}_2$, пересекающих под прямым углом ось $\vb{E}$ и образующих между собой дуальный угол $\Phi/2$.
\end{quote}
То, как мы применили дуальные аналоги к классической формуле Родрига называется принципом перенесения Котельникова-Штуди.

\section{Практическое применение}

Для описания винтового движения будем следовать принципу Котельникова-Штуди (описанного в разделе ~\ref{sec:Rodrig}) согласно которому все формулы для конечных поворотов будут иметь смысл и для винтового движения,
если в них произвести замену векторов на винты, вещественых чисел на дуальные,углы на дуальные углы и кватернионы на бикватернионы.  

\subsection{Пример №1}
\label{sec:Example}

Рассмотрим наглядный пример винтового движения ~\cite{collection:2025}.Есть некоторая прямая, которая лежит в плоскости $Oxy$ и проходит через точку $О$.Точка $O$ --- это точка начала отсчета.Она будет проходить
под углом $90^\circ$ по отношению к осям $Ox$ и $Oy$. Направляющий вектор для такой прямой будет следующим:

\begin{equation*}
    \vb{v} = ( -1, -1,  -2)^T.
\end{equation*}
Также, как выше было упомянуто у нас есть точка \(O\).Она задана следующим образом:

\begin{equation*}
    O=(0, 0, 0)^T.
\end{equation*}
Теперь вычислим момент. Мы также учитываем, что в ходе вычисления данного примера \(\vb{v}\) отложен от \(O\), из этого следует, что вычисление момента будет проходить по следующей формуле:

\begin{equation*}
    \vb{m} = \vb{p} \times \vb{v} = 
    \begin{vmatrix}
        \vb{i} & \vb{j} & \vb{k} \\
        0 & 0 & 0 \\
        -1 & -1 & -2
    \end{vmatrix}
    =
    \begin{bmatrix}
        0 \\
        0 \\
        0
    \end{bmatrix}.
\end{equation*}
Таким образом соответствующий прямой $e$ винт $e'$ можно вычислить по следующей формуле:

\begin{equation*}
    \vb{e} = \vb{v} + \varepsilon \vb{m} = ( -1, -1,  -2)^T + \varepsilon (0, 0, 0)^T = -1i -1j -2k.
\end{equation*}
Теперь рассмотрим дуальный угол для поворота прямой.Дуальный угол как выше было упомянуто представляется в виде \(A = \alpha + \varepsilon \alpha^\circ\)
На данном этапе нужно задать параметры \(\alpha\) и \(\alpha^\circ\).Где \(\alpha = \pi/2\) будет углом на который мы поворачиваем,а 
\(\alpha^\circ=1\) будет обозначать величину подъема вдоль оси.
Чтобы это реализовать,мы можем использовать формулу Родрига (вывод формулы привеен в разделе ~\ref{sec:Rodrig}):

\begin{equation*}
    \vb{e}' = \cos A \, \vb{e} + \sin A \, (\vb{E} \times \vb{e}) + 
    (1 - \cos A)(\vb{E} \cdot \vb{e}) \vb{E},
\end{equation*}
где \(\vb{E}=[0, 0, 1]\) обозначает единичный винт оси винтового перемещения.В данном случае вокруг оси Oz.
Подставим в формулу все полученные и ранее имеющиеся значения:

\begin{multline*}
    \vb{e}' = \cos A \, (-1, -1, -2) + \sin A \, ((0, 0, 1) \times (-1, -1, -2)) + (1 - \cos A)((0, 0, 1) (-1, -1, -2)) (0, 0, 1) \\
    =(0, 0, -2 + 2 \cos A) + (- \cos A, - \cos A, -2 \cos A) + (\sin A, \sin A, 0) \\
    = \begin{bmatrix}
        - \cos A + \sin A \\
        - \cos A - \sin A \\
        -2
    \end{bmatrix}.
\end{multline*}
Получив результат вычислений,мы можем воспользоваться следующими двумя формулами для упрощения вида.
\begin{quote}
    Для синуса:

    \begin{equation*}
        \sin A = \sin (\alpha + \varepsilon \alpha^\circ) = \sin \alpha + \alpha^\circ \cos \alpha \varepsilon .
    \end{equation*}
    Для косинуса:

    \begin{equation*}
        \cos A = \cos (\alpha + \varepsilon \alpha^\circ) = \cos \alpha - \alpha^\circ \sin \alpha \varepsilon .
    \end{equation*}
    \end{quote}
Используем их для разложения элементов получившегося вектор-столбца и подставим дуальный угол:

\begin{align*}
    - \cos A + \sin A = - \cos \pi /2 + \sin \pi /2 \varepsilon + \sin \pi /2 + \cos \pi /2 \varepsilon = \varepsilon + 1, \\
    - \cos A - \sin A = - \cos \pi /2 + \sin \pi /2 \varepsilon - \sin \pi /2 - \cos \pi /2 \varepsilon = \varepsilon - 1.
\end{align*}

\begin{equation*}
    \vb{e}'= 
    \begin{bmatrix}
        \varepsilon + 1 \\
        \varepsilon - 1 \\
        -2
    \end{bmatrix}
    =
    \begin{bmatrix}
        1 \\
        -1 \\
        -2 
    \end{bmatrix}
    + \varepsilon
    \begin{bmatrix}
        1 \\
        1 \\
        0
    \end{bmatrix}.
\end{equation*}
Исходя из этого,мы получаем следующее:

\begin{equation*}
    \vb{m}'=
    \begin{bmatrix}
        1 \\
        1 \\
        0
    \end{bmatrix}
    \vb{v}'=
    \begin{bmatrix}
        1 \\
        -1 \\
        -2
    \end{bmatrix}.
\end{equation*}
Теперь нормализуем изначальный вектор \(\vb{v}\):

\begin{equation*}
    \vb{v}=[-1/\sqrt{6}, -1/\sqrt{6}, -2/\sqrt{6}].
\end{equation*}
Вновь найдем \(\vb{e}'\) с учетом нормализации:

\begin{equation*}
    \vb{e}'= 
    \begin{bmatrix}
        -1/\sqrt{6} \\
        -1/\sqrt{6} \\
        -2/\sqrt{6}
    \end{bmatrix}
    + \varepsilon
    \begin{bmatrix}
        1/\sqrt{2} \\
        1/\sqrt{2} \\
        0
    \end{bmatrix}.
\end{equation*}
В полученном винте главной частью является направляющий вектор оси. Для однозначной визуализации данной оси в пространстве необходимо вычислить точку,
через которую она проходит. Это можно сделать следующим образом:
\begin{equation*}
    \vb{p}^{\prime}_{0} = \frac{\vb{v}' \times \vb{m}'}{||\vb{v}||^2} 
    =
    \begin{vmatrix}
        \vb{i} & \vb{j} & \vb{k} \\
        -1/\sqrt{6} & -1/\sqrt{6} & -2/\sqrt{6} \\
        1/\sqrt{2} & 1/\sqrt{2} & 0
    \end{vmatrix}
    =(1/3, -1/3, 1/3),
\end{equation*}
где \(\vb{p}^{\prime}_{0}\) является радиус вектором точки оси, ближайшей к началу координат O.
Полученный результат является точкой через которую будет проходить прямая \(\vb{e}'\).

\subsection{Пример №2}

Рассмотрим прямую, проходящую через точку $O = (0, 0, 1)^T$ с направляющим вектором:

\begin{equation*}
    \vb{v} = (2, 1, 0)^T.
\end{equation*}

Вычислим момент:

\begin{equation*}
    \vb{m} = \vb{p} \times \vb{v} = 
    \begin{vmatrix}
        \vb{i} & \vb{j} & \vb{k} \\
        0 & 0 & 1 \\
        2 & 1 & 0
    \end{vmatrix}
    = 
    \begin{bmatrix}
        -1 \\
        2 \\
        0
    \end{bmatrix}.
\end{equation*}

Соответствующий винт:

\begin{equation*}
    \vb{e} = \vb{v} + \varepsilon \vb{m} = (2, 1, 0)^T + \varepsilon (-1, 2, 0)^T.
\end{equation*}

Зададим дуальный угол поворота $A = \pi/2 + \varepsilon$. Применим формулу Родрига с $\vb{E} = (0, 0, 1)^T$:

\begin{multline*}
    \vb{e}' = \cos A \, (2, 1, 0) + \sin A \, 
    \begin{vmatrix}
        \vb{i} & \vb{j} & \vb{k} \\
        0 & 0 & 1 \\
        2 - \varepsilon & 1 + 2\varepsilon & 0
    \end{vmatrix}
    = -\varepsilon (2, 1, 0) + ( -1 - 2\varepsilon, 2 - \varepsilon, 0) = \\
    = \begin{bmatrix}
        -1 - 4\varepsilon \\
        2 - 2\varepsilon \\
        0
    \end{bmatrix}
    = 
    \begin{bmatrix}
        -1 \\
        2 \\
        0
    \end{bmatrix}
    + \varepsilon 
    \begin{bmatrix}
        -4 \\
        -2 \\
        0
    \end{bmatrix}.
\end{multline*}

Нормализуем исходный вектор:

\begin{equation*}
    \vb{\hat{v}} = \left( \frac{2}{\sqrt{5}}, \frac{1}{\sqrt{5}}, 0 \right)^T, \quad
    \vb{\hat{m}} = \left( -\frac{1}{\sqrt{5}}, \frac{2}{\sqrt{5}}, 0 \right)^T.
\end{equation*}

Найдем точку на преобразованной оси:

\begin{equation*}
    \vb{p}'_0 = \frac{
    \begin{vmatrix}
        \vb{i} & \vb{j} & \vb{k} \\
        -1 & 2 & 0 \\
        -4 & -2 & 0
    \end{vmatrix}
    }{5} = \frac{(0, 0, 10)^T}{5} = (0, 0, 2)^T.
\end{equation*}



\subsection{Программная реализация на Python}

Взяв за основу вычисленный ранее пример (предложенный в разделе ~\ref{sec:Example}),проведем аналогичные вычисления, но в программном виде.Для обработки дуальных чисел необходимо использовать следующий модуль ~\cite{dual_numbers}, реализующий операции над дуальными числами.


Задав начальные параметры примера, вычислим момент и винт:

\begin{minted}{python}

    m = np.cross(p, v)
    # Действительную часть берем из вектора v, а дуальную часть из момента m
    L = np.array([Dual(v[i], m[i]) for i in range(3)], dtype=Dual)
    # Теперь L — винт, задающий прямую

\end{minted}
Теперь введем дуальный угол и проведем вращение вокруг оси с помощью формулы Родрига:

\begin{minted}{python}

    def Rodrig(A, Θ, L):
    return A * (Dual(1, 0) - Dual.cos(Θ)) * np.dot(A, L) + L * Dual.cos(Θ) + np.cross(A, L) * Dual.sin(Θ)

    Θ = Dual(math.pi / 2, 1)  # дуальный угол

    # Вычисление нового винта
    L1 = Rodrig(A, Θ, L)

\end{minted}
После нормализации \(\vb{m}'\) и \(\vb{v}'\) мы можем найти точку через которую будет проходить прямая:

\begin{minted}{python}

    P = np.cross(V, M) / np.linalg.norm(V)**2

\end{minted}
Итоговый вывод программы будет выглядеть, как показано на  рисунке ~\ref{fig:img01}

\begin{figure}
    \center
    \includegraphics[width=0.5\textwidth]{img/1.pdf}
    \caption{Реализация винтового движения вокруг оси}
    \label{fig:img01}
\end{figure}

